\section{Conclusions}

ARC-ADHD is reframed as a closed-loop research framework for training and measuring recovery to goal-directed behavior after controlled challenge in children with ADHD. The core paper separates the principal regulation mechanism from four optional companion lines: facial-behavior feedback, bilateral auditory stimulation, embodied/AR context transfer, and wearable EEG measurement.

The empirical sequence is deliberately conservative. First establish that the challenge-regulation-recovery cycle is understandable and safe and that the recovery measure itself passes prespecified validity and robustness gates. Then determine feasibility and estimate design parameters. Only after the intervention, primary endpoint, active control, recovery definition, safety rules, and analysis plan are frozen should a confirmatory trial test far-transfer efficacy.

The framework is weakened by outcomes such as improvement confined to trained tasks, unstable recovery metrics, recovery change without reproducible independent transfer, external improvement without detectable change in the proposed recovery process, no clinically meaningful advantage over active control on the future primary endpoint, or unacceptable burden. These patterns are informative outcomes rather than reasons to protect the model by adding further technology. Association between change in recovery and change in an external outcome is treated as mechanistic coherence, not as proof of causal mediation.

The immediate development gates are now more specific. The project must execute the Polish EDI localization and validation pathway before inferential use; stabilize and freeze eligible-challenge and \(q_e(t)\)/\(T_{\mathrm{rec}}\) definitions without clinical-outcome-driven tuning; complete real-device BYOD benchmarking and freeze the supported Android/runtime policy; freeze the first Signal Gate task, child-facing strategy wording, and safety/overload rules; and complete the named clinical, statistical, measurement, safety, privacy, and data-governance roles required before recruitment. These gates define what must be resolved before the conceptual framework can become a participant-facing empirical study.
